% GENERATED FILE. Edit canonical lessons, then run npm run generate:books.
% canonical-source-hash: fnv1a64:6c312b25e0f0bd5f
% canonical-lessons: RU-C355-sluchaynyy, RU-C355-sluchaynost, RU-C355-sovpadenie, RU-C355-neizbezhnyy, RU-C355-voznikat

\chapter{Chance and coincidence}
\label{ch:ru-chance-and-coincidence}
\hlchaptermodality{355}

\begin{chapteropening}
\textbf{By the end of this chapter:} \emph{I can say whether something happened by chance or was bound to happen, talk about a coincidence and say that a problem arises.}

Complete the last lesson of chapter 355: I can say whether something happened by chance or was bound to happen, talk about a coincidence and say that a problem arises.
\end{chapteropening}

\section[\texorpdfstring{\ru{случайный} (slucháynyy)}{случайный (slucháynyy)}]{\ru{случайный} (slucháynyy) — accidental, chance}

\label{lesson:RU-C355-sluchaynyy}



\begin{quote}
Before the new one: say the Russian for to lower, then the Russian for to raise.
\end{quote}

\subsection*{You'll want to know: \ru{случайный}}
\textbf{\ru{случайный}} (\emph{slucháynyy}) — \textquotedblleft{}accidental, chance\textquotedblright{}.

\begin{quote}
\ru{Это} \ru{была} \ru{случайная} \ru{встреча}.
\end{quote}

\begin{quote}
\emph{Éto bylá slucháynaya vstrécha.}
\end{quote}

\begin{quote}
It was a chance meeting.
\end{quote}

\emph{\ru{Случайно}} — by chance. \emph{\ru{Ты} \ru{случайно} \ru{не} \ru{знаешь}…?} — you would not happen to know…?

\begin{cousinweb}
\textbf{\ru{случайный}} is built on \emph{\ru{случай}} \textquotedblleft{}an event, a chance\textquotedblright{}, inherited from Proto-Slavic \emph{*sŭlučajĭ}, from \emph{*lučiti}, whose sense EtymDB gives only tentatively, as \textquotedblleft{}to find, to reach\textquotedblright{} (EtymDB). A chance is what you happen to come upon.
\end{cousinweb}

\begin{grammarlens}[title={What you've built}]
One more word for telling, explaining and describing at B1.
\end{grammarlens}

\subsection*{Your turn}
\begin{itemize}
\raggedright
  \item \emph{Say it:} \emph{slucháynyy}
  \item \emph{Say it:} \emph{slucháynyy}, once more
  \item \emph{Say it:} a sentence of your own with \emph{slucháynyy}
  \item \emph{Recall:} say the Russian for to lower, then the Russian for to raise, then say \emph{slucháynyy} again
\end{itemize}

\begin{tcolorbox}[breakable,colback=teal!4,colframe=teal!35!black,title={Before you move on}]
What does \ru{случайный} mean? (\textquotedblleft{}Accidental, chance\textquotedblright{}.) Say it once more.
\end{tcolorbox}
\section[\texorpdfstring{\ru{случайность} (slucháynost)}{случайность (slucháynost)}]{\ru{случайность} (slucháynost) — chance, an accident}

\label{lesson:RU-C355-sluchaynost}



\begin{quote}
Before the new one: say the Russian for to raise, then the Russian for accidental.
\end{quote}

\subsection*{You'll want to know: \ru{случайность}}
\textbf{\ru{случайность}} (\emph{slucháynost}) — \textquotedblleft{}chance, an accident\textquotedblright{}.

\begin{quote}
\ru{Это} \ru{не} \ru{случайность}, \ru{а} \ru{план}.
\end{quote}

\begin{quote}
\emph{Éto ne slucháynost, a plan.}
\end{quote}

\begin{quote}
This is not chance but a plan.
\end{quote}

Feminine. \emph{\ru{По} \ru{случайности}} — by chance.

\begin{cousinweb}
\textbf{\ru{случайность}} is built on \emph{\ru{случайный}} with the abstract suffix \emph{-\ru{ость}}; the root is \emph{\ru{случай}}, from Proto-Slavic \emph{*sŭlučajĭ} (EtymDB).
\end{cousinweb}

\begin{grammarlens}[title={What you've built}]
One more word for telling, explaining and describing at B1.
\end{grammarlens}

\subsection*{Your turn}
\begin{itemize}
\raggedright
  \item \emph{Say it:} \emph{slucháynost}
  \item \emph{Say it:} \emph{slucháynost}, once more
  \item \emph{Say it:} a sentence of your own with \emph{slucháynost}
  \item \emph{Recall:} say the Russian for to raise, then the Russian for accidental, then say \emph{slucháynost} again
\end{itemize}

\begin{tcolorbox}[breakable,colback=teal!4,colframe=teal!35!black,title={Before you move on}]
What does \ru{случайность} mean? (\textquotedblleft{}Chance, an accident\textquotedblright{}.) Say it once more.
\end{tcolorbox}
\section[\texorpdfstring{\ru{совпадение} (sovpadéniye)}{совпадение (sovpadéniye)}]{\ru{совпадение} (sovpadéniye) — a coincidence}

\label{lesson:RU-C355-sovpadenie}



\begin{quote}
Before the new one: say the Russian for accidental, then the Russian for chance.
\end{quote}

\subsection*{You'll want to know: \ru{совпадение}}
\textbf{\ru{совпадение}} (\emph{sovpadéniye}) — \textquotedblleft{}a coincidence\textquotedblright{}.

\begin{quote}
\ru{Какое} \ru{совпадение}! \ru{Мы} \ru{родились} \ru{в} \ru{один} \ru{день}.
\end{quote}

\begin{quote}
\emph{Kakóye sovpadéniye! My rodílis v odín den.}
\end{quote}

\begin{quote}
What a coincidence! We were born on the same day.
\end{quote}

Neuter. From \emph{\ru{совпадать}}, to coincide, built on \emph{\ru{падать}}, to fall: two things that fall together.

\begin{cousinweb}
\textbf{\ru{совпадение}} is built on \emph{\ru{совпадать}} \textquotedblleft{}to coincide\textquotedblright{}, made of \emph{\ru{со}-} \textquotedblleft{}together\textquotedblright{} and \emph{\ru{падать}} \textquotedblleft{}to fall\textquotedblright{}, inherited from Proto-Slavic \emph{*padati} (EtymDB). A coincidence is two things that fall together.
\end{cousinweb}

\begin{grammarlens}[title={What you've built}]
One more word for telling, explaining and describing at B1.
\end{grammarlens}

\subsection*{Your turn}
\begin{itemize}
\raggedright
  \item \emph{Say it:} \emph{sovpadéniye}
  \item \emph{Say it:} \emph{sovpadéniye}, once more
  \item \emph{Say it:} a sentence of your own with \emph{sovpadéniye}
  \item \emph{Recall:} say the Russian for accidental, then the Russian for chance, then say \emph{sovpadéniye} again
\end{itemize}

\begin{tcolorbox}[breakable,colback=teal!4,colframe=teal!35!black,title={Before you move on}]
What does \ru{совпадение} mean? (\textquotedblleft{}A coincidence\textquotedblright{}.) Say it once more.
\end{tcolorbox}
\section[\texorpdfstring{\ru{неизбежный} (neizbézhnyy)}{неизбежный (neizbézhnyy)}]{\ru{неизбежный} (neizbézhnyy) — inevitable}

\label{lesson:RU-C355-neizbezhnyy}



\begin{quote}
Before the new one: say the Russian for chance, then the Russian for a coincidence.
\end{quote}

\subsection*{You'll want to know: \ru{неизбежный}}
\textbf{\ru{неизбежный}} (\emph{neizbézhnyy}) — \textquotedblleft{}inevitable\textquotedblright{}.

\begin{quote}
\ru{Изменения} \ru{неизбежны}.
\end{quote}

\begin{quote}
\emph{Izmenéniya neizbézhny.}
\end{quote}

\begin{quote}
Changes are inevitable.
\end{quote}

Short forms \emph{\ru{неизбежен}, \ru{неизбежна}}. Built on \emph{\ru{избежать}}, to avoid: what cannot be avoided.

\begin{cousinweb}
\textbf{\ru{неизбежный}} is built from \emph{\ru{не}-} \textquotedblleft{}not\textquotedblright{}, \emph{\ru{из}-} \textquotedblleft{}out of\textquotedblright{} and \emph{\ru{бежать}} \textquotedblleft{}to run\textquotedblright{}, inherited from Proto-Slavic \emph{*běžati}: what you cannot run away from.
\end{cousinweb}

\begin{grammarlens}[title={What you've built}]
One more word for telling, explaining and describing at B1.
\end{grammarlens}

\subsection*{Your turn}
\begin{itemize}
\raggedright
  \item \emph{Say it:} \emph{neizbézhnyy}
  \item \emph{Say it:} \emph{neizbézhnyy}, once more
  \item \emph{Say it:} a sentence of your own with \emph{neizbézhnyy}
  \item \emph{Recall:} say the Russian for chance, then the Russian for a coincidence, then say \emph{neizbézhnyy} again
\end{itemize}

\begin{tcolorbox}[breakable,colback=teal!4,colframe=teal!35!black,title={Before you move on}]
What does \ru{неизбежный} mean? (\textquotedblleft{}Inevitable\textquotedblright{}.) Say it once more.
\end{tcolorbox}
\section[\texorpdfstring{\ru{возникать} (voznikát)}{возникать (voznikát)}]{\ru{возникать} (voznikát) — to arise, to come up}

\label{lesson:RU-C355-voznikat}



\begin{quote}
Before the new one: say the Russian for a coincidence, then the Russian for inevitable.
\end{quote}

\subsection*{You'll want to know: \ru{возникать}}
\textbf{\ru{возникать}} (\emph{voznikát}) — \textquotedblleft{}to arise, to come up\textquotedblright{}.

\begin{quote}
\ru{На} \ru{работе} \ru{часто} \ru{возникают} \ru{новые} \ru{проблемы}.
\end{quote}

\begin{quote}
\emph{Na rabóte chásto voznikáyut nóvyye problémy.}
\end{quote}

\begin{quote}
New problems often come up at work.
\end{quote}

Imperfective; perfective \emph{\ru{возникнуть}} (past \emph{\ru{возник}, \ru{возникла}}). Of questions, problems and ideas.

\begin{cousinweb}
\textbf{\ru{возникать}} was taken from Church Slavonic \emph{vŭzniknǫti} (EtymDB), built from \emph{\ru{воз}-} \textquotedblleft{}up\textquotedblright{} and \emph{\ru{никнуть}} \textquotedblleft{}to bow down, to droop\textquotedblright{}: something that rises up into view.
\end{cousinweb}

\begin{grammarlens}[title={What you've built}]
One more word for telling, explaining and describing at B1.
\end{grammarlens}

\subsection*{Your turn}
\begin{itemize}
\raggedright
  \item \emph{Say it:} \emph{voznikát}
  \item \emph{Say it:} \emph{voznikát}, once more
  \item \emph{Say it:} a sentence of your own with \emph{voznikát}
  \item \emph{Recall:} say the Russian for a coincidence, then the Russian for inevitable, then say \emph{voznikát} again
\end{itemize}

\begin{tcolorbox}[breakable,colback=teal!4,colframe=teal!35!black,title={Before you move on}]
What does \ru{возникать} mean? (\textquotedblleft{}To arise, to come up\textquotedblright{}.) Say it once more.
\end{tcolorbox}
